% Felipe Muñoz Casasola
% This is free and unencumbered software released into the public domain.

% Anyone is free to copy, modify, publish, use, compile, sell, or
% distribute this software, either in source code form or as a compiled
% binary, for any purpose, commercial or non-commercial, and by any
% means.

% In jurisdictions that recognize copyright laws, the author or authors
% of this software dedicate any and all copyright interest in the
% software to the public domain. We make this dedication for the benefit
% of the public at large and to the detriment of our heirs and
% successors. We intend this dedication to be an overt act of
% relinquishment in perpetuity of all present and future rights to this
% software under copyright law.

% THE SOFTWARE IS PROVIDED "AS IS", WITHOUT WARRANTY OF ANY KIND,
% EXPRESS OR IMPLIED, INCLUDING BUT NOT LIMITED TO THE WARRANTIES OF
% MERCHANTABILITY, FITNESS FOR A PARTICULAR PURPOSE AND NONINFRINGEMENT.
% IN NO EVENT SHALL THE AUTHORS BE LIABLE FOR ANY CLAIM, DAMAGES OR
% OTHER LIABILITY, WHETHER IN AN ACTION OF CONTRACT, TORT OR OTHERWISE,
% ARISING FROM, OUT OF OR IN CONNECTION WITH THE SOFTWARE OR THE USE OR
% OTHER DEALINGS IN THE SOFTWARE.

% For more information, please refer to <https://unlicense.org/>

\documentclass{article}

\usepackage[spanish]{babel}
\usepackage[margin=3cm]{geometry}
\usepackage{microtype}
\usepackage[pdfusetitle]{hyperref}
\usepackage{minted}

\title{Copias de seguridad}
\author{Felipe Muñoz Casasola}
\date{5 de octubre de 2025; revisado el 21 de febrero de 2026}

\begin{document}
    \maketitle
    \clearpage
    \tableofcontents
    \clearpage

    % Felipe Muñoz Casasola
% This is free and unencumbered software released into the public domain.

% Anyone is free to copy, modify, publish, use, compile, sell, or
% distribute this software, either in source code form or as a compiled
% binary, for any purpose, commercial or non-commercial, and by any
% means.

% In jurisdictions that recognize copyright laws, the author or authors
% of this software dedicate any and all copyright interest in the
% software to the public domain. We make this dedication for the benefit
% of the public at large and to the detriment of our heirs and
% successors. We intend this dedication to be an overt act of
% relinquishment in perpetuity of all present and future rights to this
% software under copyright law.

% THE SOFTWARE IS PROVIDED "AS IS", WITHOUT WARRANTY OF ANY KIND,
% EXPRESS OR IMPLIED, INCLUDING BUT NOT LIMITED TO THE WARRANTIES OF
% MERCHANTABILITY, FITNESS FOR A PARTICULAR PURPOSE AND NONINFRINGEMENT.
% IN NO EVENT SHALL THE AUTHORS BE LIABLE FOR ANY CLAIM, DAMAGES OR
% OTHER LIABILITY, WHETHER IN AN ACTION OF CONTRACT, TORT OR OTHERWISE,
% ARISING FROM, OUT OF OR IN CONNECTION WITH THE SOFTWARE OR THE USE OR
% OTHER DEALINGS IN THE SOFTWARE.

% For more information, please refer to <https://unlicense.org/>
    \section{Configuración de Rclone}
    \subsection{Descripción}
    En la instalación actual de \texttt{rclone} en mi sistema, hay activadas dos sincronizaciones, una de la carpeta \texttt{\$HOME/Documentos/} y otra para el escritorio. Se sincronizan cada una a cuentas distintas de MEGA, que me ofrece los 20 Gb que necesito. Antes de subirse, todos los archivos son encriptados con una contraseña y un salteado, que se encuentran en mi base de datos de contraseñas (archivo AppImage de KeePassXC en \texttt{\$HOME/Escritorio}), que también se encuentra en el escritorio.

    \subsection{Archivos de inicio}
    Los dos archivos que inician automáticamente este sistema y lo mantienen funcionando son: \texttt{/etc/systemd/system/rclone.service} y \texttt{/etc/systemd/system/rclone\_e.service} con el siguiente contenido:

   \begin{minted}[breaklines, breakanywhere, breaksymbolleft="", breaksymbolright="", breakanywheresymbolpre="", breakanywheresymbolpost=""]{systemd}
[Unit]
After=multi-user.target
Description=rclone

[Service]
User=usuario
ExecStart=/usr/bin/rclone sync --metadata /home/usuario/Documentos/ Encriptar_documentos:
Restart=on-abort

[Install]
WantedBy=multi-user.target       
   \end{minted}
    
    Y el archivo para la sincronización del escritorio:

    \begin{minted}[breaklines=true,breaksymbolleft=""]{systemd}
[Unit]
After=multi-user.target
Description=rclone

[Service]
User=usuario
ExecStart=/usr/bin/rclone sync --metadata /home/usuario/Escritorio/ Encriptar\_escritorio:
Restart=on-abort

[Install]
WantedBy=multi-user.target
    \end{minted}

    \subsection{Uso}
    Para iniciar una interfaz web en el navegador, se puede usar el comando \texttt{rclone rcd --rc-web-gui}, en caso de pérdida, se debe intentar consultar una forma de desencriptar los archivos manualmente sin rclone, aunque por el momento no he encontrado ninguna, la forma más cómoda es restaurar el archivo de configuración.

    \subsection{Mantenimiento}
    Para actualizar la instancia, se puede ejecutar el comando \texttt{sudo rclone selfupdate}. Nótese que usa \texttt{go} para funcionar, luego no debo desinstalar nada de \texttt{go}. La configuración de \texttt{rclone} se encuentra en \texttt{\$HOME/.config/rclone/rclone.conf}, que tiene el siguiente contenido:


    
    \section{Copia de seguridad en el disco duro con \texttt{rsync}}
    En la carpeta \texttt{$\sim$/Documentos/Gestión\- del sistema/Copias de seguridad/}
    se encuentran unos guiones que crean una copia de seguridad en mi disco
    duro externo de las carpetas del escritorio, las imágenes y los documentos
    es una breve modificación de un guion sacado de la wiki de Archlinux, lo
    que hace es crear en las carpetas correspondientes copias de seguridad
    incrementales de las carpetas locales, lográndolo mediante enlaces duros
    (enlazan directamente al mismo archivo y parece que son archivos distintos)
    , luego si quiero eliminar una versión antigua, debo tener en cuenta que
    todos los archivos que no hayan sido modificados en las copias de seguridad
    siguientes también se eliminarán.

    Los guiones aparecen a continuación, primero el de los documentos, luego
    el del escritorio y luego el de las imágenes:

    \inputminted[breaklines, breakanywhere, breaksymbolleft="", breaksymbolright="", breakanywheresymbolpre="", breakanywheresymbolpost=""]
    {bash}{inc/Copia de seguridad de documentos.sh}

    \inputminted[breaklines, breakanywhere, breaksymbolleft="", breaksymbolright="", breakanywheresymbolpre="", breakanywheresymbolpost=""]
    {bash}{inc/Copia de seguridad de escritorio.sh}

    \inputminted[breaklines, breakanywhere, breaksymbolleft="", breaksymbolright="", breakanywheresymbolpre="", breakanywheresymbolpost=""]
    {bash}{inc/Copia de seguridad de imágenes.sh}

    El enlace de la wiki de Arch es \href{https://wiki.archlinux.org/title/Rsync}{este}
    que está archivado \href{https://web.archive.org/web/20251005105853/https://wiki.archlinux.org/title/Rsync}{aquí}.

    Como están guardados de forma fácil, en caso de tener que restaurarla
    es simplemente copiar los archivos que quiera.
    
    \section{Copia de seguridad completa con Rescuezilla}
    Para iniciar este proceso tengo que arrancar con el usb con una etiqueta
    que pone Rescuezilla, y una vez allí tengo gparted y se iniciará
    automáticamente una aplicación para hacer un respaldo o restaurar un
    disco duro.
\end{document}
